(1) 氢样离子的点核Dirac束缚能谱给为
\[
E_{nj}=mc^2\left[1+\frac{(Z\alpha)^2}{\bigl(n-j-\tfrac12+\sqrt{(j+\tfrac12)^2-(Z\alpha)^2}\bigr)^2}\right]^{-1/2},
\]
其中 \(n\ge1,\ j=1/2,3/2,\ldots,n-1/2\)，\(Z\alpha\ll1\)。系统地展开至 \((Z\alpha)^4\)，分别识别非相对论主项与精细结构项。
(2) 对 \(n=2\) 的 \(j=1/2\) 和 \(j=3/2\) 两组态求精细结构能量差的最低非零阶，证明哪组能量更高，说明其与 \(m_j\) 简并的关系。
(3) 计算相对论精细结构与非相对论束缚能的量级比，讨论 \(Z\to0\) 的形式极限及 \(n\) 固定时核电荷增强的标度；解释为何同 \(n,j\) 不同轨道角动量的态在理想Dirac点核模型中仍可能简并。
(4) 从谱中平方根的实数条件给出最低 \(j=1/2\) 通道的点核临界条件，说明高 \(Z\) 下有限核尺寸、真空极化与Lamb位移为何使此公式不能无限外推。
